\section*{Chapter \ref{ch:g3:stones-pebbles-sand} --- Stones, Pebbles and Sand}

\begin{solution}{exo:g3:stones-pebbles-sand:1}
The \omterm{def:g3:stones-pebbles-sand:rock}{bedrock} is the solid \omterm{def:g3:stones-pebbles-sand:rock}{rock} that lies under the soil everywhere. We see
it where there is no soil to hide it: on the steep sides of \omterm{def:g1:landscapes-around-us:landform}{mountains}, in
\omterm{def:g2:water-on-earth:ocean}{sea} cliffs, along roads cut through a \omterm{def:g1:landscapes-around-us:landform}{hill}.
\end{solution}

\begin{solution}{exo:g3:stones-pebbles-sand:2}
\omterm{def:g3:stones-pebbles-sand:sand}{Mud}, \omterm{def:g3:stones-pebbles-sand:sand}{sand}, \omterm{def:g3:stones-pebbles-sand:pebble}{pebble}, boulder.
\end{solution}

\begin{solution}{exo:g3:stones-pebbles-sand:3}
1 millimetre is smaller than 2 millimetres: \omterm{def:g3:stones-pebbles-sand:sand}{sand}. 3 millimetres is bigger
than 2 millimetres: \omterm{def:g3:stones-pebbles-sand:pebble}{gravel}.
\end{solution}

\begin{solution}{exo:g3:stones-pebbles-sand:4}
Natural: the \omterm{def:g3:stones-pebbles-sand:pebble}{pebble}, the block of granite (cut by people, but the \omterm{def:g3:stones-pebbles-sand:rock}{rock} is
natural), the \omterm{def:g3:stones-pebbles-sand:sand}{sand} on the beach. Made by people: the brick and the glass.
\end{solution}

\begin{solution}{exo:g3:stones-pebbles-sand:5}
The water of the \omterm{def:g1:landscapes-around-us:water}{river} has rolled it along the bottom and knocked it
against other stones for a long time: its corners have worn away.
\end{solution}

\begin{solution}{exo:g3:stones-pebbles-sand:6}
Granite, for a wall; thin, flat stone slabs (a \omterm{def:g3:stones-pebbles-sand:rock}{rock} like slate), for roofs.
\end{solution}

\begin{solution}{exo:g3:stones-pebbles-sand:7}
The round one: the longer a stone is rolled by water, the more its corners
wear away. The sharp one has not travelled far from the \omterm{def:g3:stones-pebbles-sand:rock}{rock} it broke from.
\end{solution}

\begin{solution}{exo:g3:stones-pebbles-sand:8}
The biggest \omterm{def:g3:stones-pebbles-sand:pebble}{pebble} is 3 centimetres long. One centimetre is 10
millimetres, and $10 \div 2 = 5$: five grains of 2 millimetres fit along
1 centimetre.
\end{solution}

\begin{solution}{exo:g3:stones-pebbles-sand:9}
In the quiet pond. Fast water keeps the tiny grains of \omterm{def:g3:stones-pebbles-sand:sand}{mud} moving; only
where the water is almost still can they settle to the bottom.
\end{solution}

\begin{solution}{exo:g3:stones-pebbles-sand:10}
Because they are made from materials taken from the ground: bricks from
clay, concrete from \omterm{def:g3:stones-pebbles-sand:sand}{sand}, \omterm{def:g3:stones-pebbles-sand:pebble}{gravel} and cement, glass from \omterm{def:g3:stones-pebbles-sand:sand}{sand}. Clay, \omterm{def:g3:stones-pebbles-sand:sand}{sand} and
\omterm{def:g3:stones-pebbles-sand:pebble}{gravel} come from broken \omterm{def:g3:stones-pebbles-sand:rock}{rock}.
\end{solution}
